% !TeX root = ../main.tex
\section{CacheWise Replication: Experiments and Results}
\label{sec:replication}

% Evidence and figure-generation procedures are recorded in notes/README.md.
% The primary serving results are the consolidated six-policy study, not job 10807.
\subsection{Study Scope and Workload Characterization}
We study whether the scheduling and eviction mechanisms of CacheWise improve serving efficiency in an adapted coding-agent replay.
Our evaluation asks how much benefit comes from prefix-aware scheduling, whether predictive eviction adds to that benefit, and how recorded tool-duration information and queue aging affect the outcome.
The experiment reproduces these mechanisms and comparisons under a different hardware and model configuration; it does not attempt to recreate the original paper's numerical results exactly.

The local trace release contains 94 source files and 20,634 LLM event records representing 11,453 distinct messages.
For descriptive analysis, we group records by source trace and message identifier, discard synthetic-model messages, and retain messages with positive prompt and output usage.
Prompt length includes input, cache-creation, and cache-read tokens; we take the maximum observed usage within each message.
This yields 11,388 messages from 80 source traces.
The descriptive population includes both training and held-out traces and is separate from the serving evaluation cohort.

Figure~\ref{fig:replication-workload} shows a strong difference between input and output lengths: the median prompt contains 80,670 tokens, whereas the median output contains 26.
The median trace with included messages contains 98.5 unique messages.
These properties make prefix reuse relevant: processing a long input afresh at each turn can introduce substantial work even when the generated response is short.
The individual tool-duration distributions also differ across common tools, providing a motivation for using tool information when deciding which state to retain.
These descriptive observations do not establish that a particular predictor improves serving performance.

\begin{figure*}[t]
  \centering
  \includegraphics[width=\textwidth]{replication-workload.pdf}
  \caption{Workload characterization regenerated from the local trace release.
  Panels (a) and (b) use 11,388 deduplicated, non-synthetic messages with positive usage from 80 of the 94 source traces.
  Panel (c) shows the six most frequent tools with positive completed durations associated with those messages; repeated tool identifiers are merged, retaining the earliest launch for the same result.
  Boxes span the 25th--75th percentiles and whiskers the 5th--95th percentiles; outliers are omitted.
  These corpus-level plots are descriptive and do not represent the smaller serving cohort.}
  \label{fig:replication-workload}
\end{figure*}

\subsection{Experimental Setup and Replay Procedure}
The primary serving experiment uses ten complete held-out segments drawn from eight source sessions.
Each run contains 83 requests, 4,894,853 prompt tokens counted across requests, and 15,235 generated tokens.
Repeated prefixes contribute repeatedly to the prompt-token total; that quantity is not the number of distinct input tokens.
The selected workload has 73 between-turn transitions, 57 of which have recorded delays of at most one second.
It therefore emphasizes short reuse intervals, and is not assumed to represent all tool-delay distributions in the release.

We use Qwen2.5-Coder-14B-Instruct in BF16 on one RTX Pro 6000 Blackwell Max-Q GPU per repeat, with tensor parallelism of one.
Prefix caching and 512-token chunked prefill are enabled.
The KV budget is 17,981 blocks of 16 tokens, equivalent to 287,696 tokens.
The ten initial prompts total 487,696 tokens, approximately 1.70 times this capacity.
Thus, $N=10$ already creates memory pressure on this testbed; equal session counts do not imply equal contention across different hardware configurations.
Table~\ref{tab:replication-setup} summarizes the principal differences from the reference experiment.

\begin{table}[t]
  \centering
  \small
  \setlength{\tabcolsep}{3pt}
  \hyphenpenalty=10000
  \begin{tabular}{>{\raggedright\arraybackslash}p{0.23\columnwidth}>{\raggedright\arraybackslash}p{0.32\columnwidth}>{\raggedright\arraybackslash}p{0.32\columnwidth}}
    \toprule
    Setting & Reference & This study \\
    \midrule
    Accelerator & Two H200 GPUs & One RTX Pro 6000 per repeat \\
    Model size & 32B & 14B \\
    Tensor parallelism & Two & One \\
    Prefill chunk & 512 tokens & 512 tokens \\
    Predictor inputs & Tool names and arguments & Tool names; arguments redacted \\
    Training membership & 80\% of sessions & 34 of 94 source traces \\
    Serving loads & Multiple values of $N$ & One $N=10$ cohort \\
    \bottomrule
  \end{tabular}
  \caption{Selected configuration differences from CacheWise~\cite{Tiwari2026CacheWise}.
  Exact membership of the reference serving split is unavailable; the local training fraction uses the 94 available source traces as its denominator.}
  \label{tab:replication-setup}
\end{table}

All policies run in the same CacheWise engine fork, pinned to revision \texttt{a4b3affa} with recorded local patches.
The patches address session lifecycle handling, protection of queued returning sessions, and the timestamp used to observe tool durations.
Each policy starts a fresh server, receives the same selected workload and token lengths, preserves external delays, and marks completed segments idle.
The baseline is therefore an FCFS/LRU control within this engine, rather than a comparison against a separate stock-vLLM installation.

Because the released text is redacted, the replay uses seeded synthetic prompts that preserve within-segment prefix growth.
Generation is constrained to fixed token identifiers and recorded output lengths.
Turns within a segment run sequentially, while the ten segments execute concurrently.
Actual human-input boundaries and context shrink split the traces into segments, and overlapping tools contribute the interval from their earliest launch to their latest result.
This procedure preserves a length-and-reuse workload, rather than original model answer quality.
The predictor bootstrap uses 34 training source sessions that are excluded from serving evaluation; the learned predictor uses tool names because argument content is unavailable.

We evaluate six settings, each with three repeats, for 18 full runs and 1,494 request executions.
The four core settings form a two-by-two comparison of FCFS versus prefix-aware scheduling and LRU versus learned predictive eviction.
The remaining settings replace learned predictions with recorded tool durations or enable the local 30-second queue-aging override.
Within each repeat, all six policies share one physical GPU; the three repeat blocks run on separate GPUs with matching model, memory, driver, runtime, engine, and replay configuration.
Policy order is partially rotated, not fully balanced.

\subsection{Metrics and Evidence Validation}
Our primary metric is mean session serving time: the mean, across segments, of the sum of their LLM request latencies, averaged over three repeats.
It excludes external tool execution.
Whole-workload elapsed time includes those waits and measures a different outcome.
Prompt-token hit rate is computed from completed-request usage as cached prompt tokens divided by total prompt tokens, rather than from cache-lookup counters.

For supporting metrics, useful-token goodput counts each new prompt token and generated output token once per segment and divides by whole-workload elapsed time.
The selected workload contains 605,956 distinct prompt tokens by this accounting and 621,191 useful tokens including output.
Latency percentiles use linear interpolation within each 83-request run; reported summaries average the three per-run percentiles.
They are not percentiles of averaged request latencies, and observed maxima are reported separately.

The experiment audit verifies identical request hashes and lengths, 6,811 successful metric snapshots, zero failed snapshots, and zero preemptions.
When preparing the manuscript, we also recalculated the per-run means, token totals, and cache-hit rates from the original request summaries.
These checks validate consistency of the recorded experiment.
They do not establish that the engine is equivalent to the original paper's executable.
All repeat ranges below are observed minima and maxima, not confidence intervals; the matched-repeat plots preserve the policy comparisons within each GPU block.

\subsection{Main Results and Scheduling--Eviction Ablation}
\begin{table*}[t]
  \centering
  \small
  \begin{tabular}{lllrrr}
    \toprule
    Scheduler & Eviction & Queue aging & Mean (s) & Repeat range (s) & Hit rate (\%) \\
    \midrule
    FCFS & LRU & --- & 595.04 & 589.87--599.00 & 64.689 \\
    FCFS & Learned predictive & --- & 529.97 & 523.43--534.05 & 69.759 \\
    Prefix-aware & LRU & Disabled & \textbf{337.87} & 332.97--342.87 & 79.009 \\
    Prefix-aware & Learned predictive & Disabled & 349.67 & 343.52--357.31 & 79.048 \\
    Prefix-aware & Tool-duration oracle & Disabled & 349.58 & 344.42--357.05 & 79.016 \\
    Prefix-aware & Learned predictive & 30 seconds & 599.14 & 592.46--604.94 & 64.984 \\
    \bottomrule
  \end{tabular}
  \caption{Consolidated serving results at $N=10$, with three repeats per setting.
  Mean and range refer to mean session serving time, excluding tool execution.
  Queue aging overrides prefix priority for sufficiently old waiting requests; it is not cache expiration.
  All comparisons use the same engine and workload within each repeat.}
  \label{tab:replication-results}
\end{table*}

Table~\ref{tab:replication-results} and Figure~\ref{fig:replication-serving} show that prefix-aware scheduling provides the largest reduction in mean session serving time in this cohort.
With LRU eviction unchanged, replacing FCFS scheduling reduces the mean from 595.04 to 337.87 seconds, a $1.761\times$ time ratio, or a 43.22\% reduction.
This comparison isolates the scheduling choice while retaining the same eviction policy.

The effect of predictive eviction depends on that scheduling choice.
Under FCFS, replacing LRU with learned predictive eviction reduces the mean to 529.97 seconds, a 10.94\% improvement.
Under prefix-aware scheduling, the same replacement increases the mean from 337.87 to 349.67 seconds, a 3.49\% regression.
Both directions hold in all three matched repeats.
Learned CacheWise without aging still improves over FCFS/LRU by a $1.702\times$ time ratio, but does not improve the primary metric over the prefix-aware LRU control.
The results therefore demonstrate an interaction between scheduling and eviction in the tested workload, rather than independent additive gains from both mechanisms.

\begin{figure*}[t]
  \centering
  \includegraphics[width=\textwidth]{replication-serving.pdf}
  \caption{Serving comparison from the consolidated six-policy experiment.
  Panel (a) shows mean session serving time with whiskers spanning the three observed repeat means.
  Panel (b) shows percentage changes within each matched GPU repeat; negative values indicate lower time than the corresponding control.
  Predictive eviction helps FCFS but increases mean time under prefix-aware scheduling in every repeat.
  The tool-duration oracle shows no consistent advantage over the learned predictor.
  Whiskers are observed ranges, not confidence intervals.}
  \label{fig:replication-serving}
\end{figure*}

The cache-hit metric alone does not explain the session-time ordering.
Learned CacheWise without aging has a slightly higher hit rate than prefix-aware LRU, 79.048\% versus 79.009\%, yet a higher mean session time.
Whole-workload elapsed time gives the opposite ordering, averaging 642.91 seconds for learned CacheWise and 646.28 seconds for prefix-aware LRU.
Corresponding mean useful-token goodputs are 966.58 and 961.32 tokens/s.
These small elapsed-time and goodput differences should not be substituted for the primary mean-session result.
They show that finishing the entire batch and minimizing the average segment's accumulated serving time are distinct objectives.

\subsection{Oracle and Queue-Aging Diagnostics}
Replacing learned predictions with recorded tool durations yields a mean of 349.58 seconds, only 0.028\% below the learned predictor's unrounded mean.
The oracle is slower in one repeat and faster in two.
Thus, providing recorded tool durations does not yield a consistent material advantage in this cohort.
The oracle supplies tool-duration information, not exact future server-arrival times or an optimal eviction sequence; the result does not establish that prediction is irrelevant in other workloads.

The local queue-aging rule has a larger effect.
With the learned predictor fixed, disabling the 30-second override reduces mean session serving time from 599.14 to 349.67 seconds, a 41.64\% reduction.
Average per-run P99 request latency also decreases, from 268.19 to 233.03 seconds.
Nevertheless, the largest individual request latency across repeats increases from 336.92 to 466.45 seconds.
The maximum and P99 therefore support different statements about tail behavior.
This experiment diagnoses a responsiveness tradeoff in the executed queue implementation; the published description does not identify this particular override as part of the evaluated design.

\subsection{Interpretation, Limitations, and Transition}
The replication supports a benefit from prefix-aware scheduling on this memory-constrained workload, while the added benefit of predictive eviction is conditional on the scheduler.
It does not reproduce the reference paper's additional predictive-eviction benefit over prefix-aware LRU~\cite{Tiwari2026CacheWise}.
This is a result for the adapted setup, not a refutation of the original evaluation.
Hardware, model size, training membership, argument redaction, engine patches, and cohort composition differ.
Only one $N=10$ cohort is measured, so higher-concurrency results remain untested.

There is also a concrete predictor-preparation limitation: bootstrap samples describe individual tools, whereas runtime metadata can describe parallel batches with repeated tool names.
A prior input audit found 16 of the selected cohort's 73 transitions lacked a matching initial bootstrap key.
Its isolated performance impact has not been measured, and argument-level prediction cannot be recovered by clustering identical redacted placeholders.
Broader load and cohort coverage, matched training/runtime keys, and complete eviction instrumentation would allow stronger conclusions about when predictive eviction helps.

Finally, substantial uncached prefill work remains under every policy.
Learned CacheWise without aging computes approximately 1.03 million prompt KV tokens per run, including first-time prefill.
This is not a full eviction count or a measurement of transferred bytes; repeated prefill and block-alignment or last-token handling require separate accounting.
No exclusive CPU/GPU cost profile or completed DRAM/SSD offloading comparison is available.
These limits motivate the recovery-cost question in Section~\ref{sec:tradeoff}, while leaving its numerical answer to the subsequent study.
